\section{Related Work}

\subsection{Automated Scientific Research}
Recent breakthroughs in LLMs \citep{reasoning-survey,long-cot-survey} have propelled them into a central position within the domain of Automated Scientific Discovery \citep{dsgym,agent-laboratory,ai-scientist,how-far-ai-sci,datamind}.
The complete workflow of automated scientific discovery comprises five consecutive phases: \textit{i) Literature Reviewing}, during which LLMs search for papers on designated topics across the internet or specialized collections and consolidate them into organized surveys \citep{autosurvey,surveyx,opensholar,surveyforge,litllms}; \textit{ii) Hypothesis Generation}, where LLMs leverage both their inherent parametric knowledge and acquired external information to formulate feasible research concepts \citep{chain-of-ideas,virsci,deep-ideation,researchagent,scipip}; \textit{iii) Method Implementation}, wherein LLMs convert the generated hypotheses into functional code, verify them via rigorous experimentation, and conduct statistical evaluation \citep{alphaevolve,automind,ml-master,alpharesearch,aide}; \textit{iv) Manuscript Writing}, in which LLMs document the research rationale, technical approach, and experimental outcomes in the form of academic papers or reports \citep{overleafcoplilot,xtragpt}; and \textit{v) Peer Reviewing}, where LLMs assume the responsibilities of reviewers to assess manuscripts from multiple perspectives \citep{cycleresearcher,agentreview,reviewer2,deepreview}.
The entire workflow of automated scientific discovery is an extremely knowledge-intensive process, in which literature review serves as the primary source of external knowledge beyond model parametric knowledge.
Consequently, a precise scientific search is of paramount importance for the whole workflow.

\subsection{Scientific Search and Discovery}
Human scientists typically conduct scientific retrieval through general-purpose academic search platforms such as Google Scholar and Semantic Scholar, domain-specific preprint servers including arXiv, ChemRxiv, and PubMed, or official publisher platforms for journals and conferences.
In the domain of automated scientific research, early efforts primarily relied on keyword or vector-based retrieval within local paper collections \citep{researchagent,virsci,rnd,scipip}.
With the agentic advancement of LLMs, web-based literature resources have become accessible through API calling \citep{can-llm-gen-novel,chain-of-ideas,deep-ideation,ai-researcher,internagent,innoeval}.
Deep research agent frameworks can further leverage the semantic understanding and reasoning capabilities of LLMs to enable in-depth literature retrieval \citep{deepxiv,wispaper,opennovelty}.
However, these approaches not only incur high computational costs and response latency but also, due to the absence of deterministic cognitive maps as anchors for LLMs, render them highly susceptible to logical hallucinations within complex exploratory trajectories.
So we argue that KG is an indispensable knowledge organization form for scientific discovery because, although LLMs have achieved remarkable advancements in semantic understanding, they still exhibit substantial deficiencies in capturing logical relationships among knowledge entities, a capability of paramount importance for scientific research that transcends mere semantic associations.
A recent related work, OmniScientist \citep{omniscientist}, has also proposed a research knowledge base.
However, it lacks the integration of core keywords for paper interconnection and semantic vectors.
Furthermore, its Elasticsearch-based search algorithm merely relies on simple propagation through citation and reference relationships, without performing structured traversal and deep topological reasoning over heterogeneous subgraphs to uncover potentially relevant literature.